\section{Princípio de Decomposição}

\begin{frame}
\frametitle{O Conceito de Decomposição}
Quando o problema excede a capacidade de memória até do Simplex Revisado, usamos a \textbf{Decomposição de Dantzig-Wolfe}.
\vspace{0.3cm}
\begin{itemize}
    \item \textbf{Cenário:} Grandes corporações descentralizadas.
    \item \textbf{Ideia Central:} Em vez de resolver tudo de uma vez, o problema é "quebrado" em pedaços menores que negociam entre si.
    \item \textbf{Estrutura da Matriz:} Bloco-Angular.
\end{itemize}
\end{frame}

\begin{frame}
\frametitle{A Estrutura Bloco-Angular}
Imagine uma empresa com várias fábricas. Cada uma tem suas regras ($A_k$), mas compartilham um orçamento global ($L$).

\[
A = \begin{bmatrix}
L_1 & L_2 & \dots & L_K \\ \hline
A_1 & 0 & \dots & 0 \\
0 & A_2 & \dots & 0 \\
\vdots & \vdots & \ddots & \vdots \\
0 & 0 & \dots & A_K
\end{bmatrix}
\]
\end{frame}

\begin{frame}
\frametitle{Fundamento Matemático: Convexidade}
O método se baseia no \textbf{Teorema da Representação}:
\vspace{0.2cm}
\begin{block}{Vértices do Poliedro}
    Qualquer solução viável dentro de um bloco $k$ pode ser escrita como uma \textbf{combinação convexa} dos vértices ($v$) desse bloco.
    \[ x_k = \sum \lambda_{j} v_{kj} \quad \text{onde} \quad \sum \lambda = 1 \]
\end{block}
\vspace{0.2cm}
O "Problema Mestre" passa a procurar os pesos perfeitos ($\lambda$) para misturar esses vértices, em vez de procurar $x$ diretamente.
\end{frame}

\begin{frame}
\frametitle{O Algoritmo de Decomposição (Ciclo)}
O processo imita um mercado econômico interno:

\begin{enumerate}
    \item \textbf{Mestre Restrito (Sede):} Resolve o problema atual com os vértices que já conhece. Gera preços sombra ($\pi$) para os recursos globais.
    
    \item \textbf{Subproblemas (Fábricas):} Recebem os preços $\pi$. Otimizam internamente para minimizar seu custo ajustado:
    \[ \min (c_k - \pi^T L_k)x_k \]
    
    \item \textbf{Geração de Colunas:} Se uma fábrica encontra uma solução que melhora o custo global, essa solução (novo vértice) é enviada ao Mestre.
    
    \item \textbf{Iteração:} O ciclo se repete até nenhuma fábrica conseguir melhorar o preço.
\end{enumerate}
\end{frame}

\section{Conclusão}

\begin{frame}
\frametitle{Conclusão}
\begin{itemize}
    \item O \textbf{Simplex Revisado} é uma evolução necessária para a computação moderna, focando na álgebra matricial esparsa e sob demanda.
    \item A \textbf{Decomposição de Dantzig-Wolfe} permite resolver problemas de escala massiva explorando a estrutura geométrica dos poliedros (convexidade).
    \item Juntos, formam a base dos solvers comerciais (Gurobi, CPLEX) utilizados em logística, energia e planejamento industrial.
\end{itemize}
\end{frame}